\section{서론}

카메라 기반 인간 센싱 시스템은 널리 활용되고 있지만, 사생활 침해 우려,
조명 변화, 시야 가림 등의 문제로 실제 환경에서의 활용이 제한될 수 있다.
이에 따라 Wi-Fi 센싱은 비접촉·비시각 방식의 대안으로 주목받고 있다.
Wi-Fi 센싱은 사람의 움직임으로 인해 발생하는 무선 신호 전파 변화를
활용하므로, 사용자가 별도의 센서를 착용하지 않아도 기존 통신 인프라를
이용하여 다양한 인간 센싱 응용을 지원할 수 있다.

대부분의 Wi-Fi 센싱 연구는 채널 상태 정보(Channel State Information,
CSI)를 사용한다. CSI는 수신 신호 세기 지표(Received Signal Strength
Indicator, RSSI)와 같은 거친 측정값과 달리 송수신 안테나 쌍에 대한
부반송파 단위의 진폭 및 위상 응답을 제공한다 \cite{halperin2011tool}.
이러한 측정값의 변화는 실내 위치 추정 \cite{kotaru2015spotfi}, 호흡
모니터링 \cite{wang2017tensorbeat}, 제스처 인식 \cite{zheng2019widar3} 및
인간 행동 인식 \cite{wang2015understanding} 등 다양한 센싱 응용에 활용되어
왔다.

CSI는 높은 센싱 성능을 제공하지만, 다수의 상용 Wi-Fi 장치에서 쉽게
획득할 수 있는 정보는 아니다. CSI 추출에는 일반적으로 지원되는 네트워크
인터페이스 카드, 수정된 펌웨어, 모니터 모드 설정 또는 전용 CSI 수집 도구가
필요하다 \cite{gringoli2019free}. 이러한 요구 사항은 CSI 기반 센싱
시스템의 실제 배치를 제한할 수 있다. 최근에는 이에 대한 대안으로 빔포밍
피드백 정보(Beamforming Feedback Information, BFI)를 활용하는 연구가
진행되고 있다 \cite{beamsense2023,bfmsense2024,beamsense2025,wibfi2023}.
IEEE 802.11ac/ax 네트워크에서 단말은 sounding 신호를 이용하여 무선
채널을 추정하고, 압축된 빔포밍 피드백을 액세스 포인트에 전송한다
\cite{ieee2021ax,ieee2013ac}. BFI는 표준화된 빔포밍 절차에서 생성되기
때문에, Wi-Fi 통신 과정에서 이미 교환되는 정보를 센싱에 활용할 수 있다는
장점이 있다.

BFI는 무선 채널 정보를 압축된 형태로 표현한다. 피드백 오버헤드를 줄이기
위해 빔포밍 행렬은 일반적으로 $\phi$와 $\psi$로 표현되는 양자화 각도
파라미터를 사용하여 부호화된다 \cite{itahara2022beamforming,li2024efficient}.
이러한 빔포밍 피드백 각도(Beamforming Feedback Angle, BFA)는 복소수
빔포밍 행렬 $\mathbf{V}$를 나타낸다. $\mathbf{V}$는 추정된 무선 채널의
우측 특이 행렬에 해당하며, 빔포밍에 사용되는 송신 측 공간 방향을 표현한다.
이 압축 표현에는 사람의 행동에 영향을 받는 채널 의존 정보가 남아 있지만,
데이터 구조와 물리적 의미는 기존 CSI와 다르다. 따라서 BFI는 빔포밍 표현의
특성을 고려하여 처리해야 한다.

다른 데이터 기반 무선 센싱 방식과 마찬가지로 BFI 기반 행동 인식의 성능도
충분한 라벨 학습 데이터의 확보 여부에 크게 좌우된다. 특히 다양한 행동,
참가자 및 배치 환경을 포함하는 무선 측정 데이터를 수집하고 라벨링하는
과정에는 많은 시간이 필요하다. 잡음 추가, 스케일 조정 또는 원소 단위 혼합과
같은 기존 증강 기법은 샘플 수를 늘릴 수 있지만, BFA의 순환 양자화 구조와
시간적 변화를 명시적으로 모델링하지 않는다. 따라서 downstream BFI 센싱
모델이 사용하는 데이터 구조에 특화된 생성 기법이 필요하다.

RF-Diffusion은 RF 시계열 신호의 시간·주파수 영역 및 복소수 특성을
학습함으로써 생성형 데이터 증강의 가능성을 보였다 \cite{chi2024rfdiffusion}.
본 연구에서는 비교적 접근성이 높은 BFI로 이 방식을 확장하기 위해, 먼저
복원된 복소수 빔포밍 행렬 $\mathbf{V}$를 생성 대상으로 사용하는 방법을
검토하였다. 그러나 두 데이터가 모두 복소수라는 사실만으로 $\mathbf{V}$와
CSI가 구조적으로 동일한 것은 아니다. CSI는 채널 응답 자체를 나타내는 반면,
$\mathbf{V}$는 채널의 공간 방향을 나타내며 정규화 및 기하학적 제약과 위상
모호성을 갖는다. 초기 실험에서 복소 행렬을 직접 생성한 결과, 예상된 텐서
형태와 기본적인 수치 유효성은 유지되었지만 BFI 센싱에 필요한 시간적 beam
변화와 행동 의미를 안정적으로 보존하지 못했다. 이는 입력 형식만 변경하여
CSI 중심의 RF 생성 모델을 BFI에 적용하는 방식으로는 충분하지 않음을
보여준다.

이 문제를 해결하기 위해 본 연구에서는 BFA 시퀀스 생성에 특화된 센싱 인지형
디퓨전 프레임워크인 \textbf{BeamDiff}를 제안한다. BeamDiff는 절대 BFA
프레임을 각각 독립적으로 생성하지 않는다. 대신 10프레임 window의 첫 실제
프레임을 anchor로 유지하고, 이후 9개의 프레임 간 변화량을 생성한다. 각
변화량은 BFA 채널별 양자화 범위를 고려한 최단 경로 순환 delta로 표현하여
각도 경계에서 발생하는 인위적인 불연속을 방지한다. Denoising network는
첫 프레임 anchor, diffusion timestep 및 목표 행동 조건을 함께 입력받는다.
또한 표준 diffusion noise-prediction objective에 clean-delta reconstruction
loss, temporal-fidelity loss 및 frozen BeamSense classifier의 supervision을
결합한다. 생성된 변화량을 anchor부터 순차적으로 누적한 뒤 유효한 양자화
BFA 범위로 변환함으로써, 기존 BFI 센싱 모델에 직접 입력할 수 있는 합성
샘플을 생성한다.

BeamDiff는 CSI-BFI-HAR 데이터셋의 HAR-1 및 HAR-3 환경에서 평가하였다.
모든 실험은 real-only와 real-plus-synthetic 조건을 paired protocol로
비교했으며, 실제 validation 및 test split은 동일하게 고정하고 train split에만
합성 데이터를 추가하였다. HAR-1의 1:1 증강에서는 5개의 모델 초기화 seed에
대한 BeamSense 평균 정확도가 91.62\%에서 94.63\%로 향상되어 평균
3.00\%p의 paired gain을 보였으며, 5개 seed 중 4개에서 성능이 향상되었다.
HAR-3의 초기 paired 평가에서는 정확도가 98.70\%에서 99.65\%로 향상되고
test error가 81개에서 22개로 감소하였다. 이러한 결과는 BFA 특화 생성
방식이 동일 환경 내 센싱 성능을 높일 수 있음을 보여준다. 한편 참가자가 다른
조건의 실험에서는 성능이 크게 감소하여, 미관측 참가자에 대한 일반화는 데이터
증강과 구분하여 다루어야 할 더 어려운 문제임을 확인하였다.

본 연구의 주요 기여는 다음과 같다.

\begin{itemize}
    \item 복소 RF 생성형 증강 기법을 BFI로 확장하는 과정을 분석하고,
    일반적인 복소 행렬 생성과 BFI 센싱에 필요한 beam 방향 변화 사이의
    구조적 불일치를 확인하였다.

    \item 양자화된 프레임을 독립적으로 생성하는 대신, 실제 첫 프레임을
    anchor로 사용하여 최단 경로 순환 각도 변화량을 생성하는 BFA 특화 조건부
    디퓨전 프레임워크 BeamDiff를 제안한다.

    \item Diffusion noise prediction, clean-delta reconstruction,
    temporal-fidelity regularization 및 frozen classifier supervision을
    결합하여 BFA의 시간적 특성과 행동 의미를 함께 보존하는 센싱 인지형
    학습 목적 함수를 설계한다.

    \item 두 BFI 센싱 환경에서 paired downstream evaluation을 수행하고,
    증강 이득, 생성 데이터 fidelity, 모델 seed에 따른 변화 및
    cross-participant 조건의 한계를 분석한다.
\end{itemize}

% 최종 제출 전 HAR-3 단일 seed 결과를 다중 seed 평균과 분산으로 교체하고,
% 모든 BibTeX key가 실제 참고문헌 파일과 일치하는지 확인해야 한다.
